\section{Empirical Strategy}
\label{sec:empirical}

The analysis has two stages. We first estimate the urgent-vs-ordinary
cost margin under item, time, and buyer fixed effects. We then bound
the under-the-gun (UTG) gap between litigated and administrative urgent
purchases under selection into the administrative channel. The
identification of the within-item gap is straightforward; the
identification of the UTG gap requires explicit selection bounds.

\subsection{Identification: urgent vs.\ ordinary}

Identification of the urgent-vs-ordinary premium requires that,
conditional on item, time, and PBU fixed effects, urgency is
uncorrelated with unobserved determinants of procurement outcomes.
Three institutional features support this. First, court orders and
administrative requests originate from patients' health needs and legal
or scientific-committee assessments---not from the procurement process.
Second, urgent demand is poorly correlated with PBU-level mismanagement:
only \BPstockShortageShare{} of claims arise from stock shortages or
service failures. Third, the urgency regime is determined externally
(by judges or the SES/SP scientific committee), not by the purchasing
unit. We frame this design as a within-item descriptive estimate of the
cost margin associated with the urgent regime.

A Borusyak-Jaravel-Spiess (\citeyear{borusyak2024revisiting}) imputation
event study around each item's first court order provides a dynamic
check (Online Appendix Figure~\ref{fig:event_study}). Pre-period
coefficients are economically small (max absolute pre-period coefficient
$=\BPbjsPreMaxAbs$); post-period coefficients rise from $\BPbjsETzero$
at $t=0$ to $\BPbjsETfive$ at $t=+5$. We honest-test this dynamic
pattern using Rambachan-Roth (\citeyear{rambachan2023more}) sensitivity:
the breakdown linear-extrapolation parameter at which the $t=0$
confidence interval just contains zero is $M=\BPrrBreakdownM$. The
$t=0$ effect is significantly positive at zero pre-period violation
but does not survive a linear extrapolation of the observed pre-period
maximum. We honor this by treating the cross-sectional fixed-effects
estimate as the primary object and the BJS event-study as supportive
dynamic evidence rather than a stand-alone identification claim.

\subsection{Identification: under-the-gun}

Both litigated and administrative urgent purchases face identical
planning constraints---compressed timelines, small quantities, diverted
budget resources, and the same BEC platform and bidding procedures.
They differ only in penalty exposure. The within-item UTG comparison
isolates the cost margin associated with the sanction regime, conditional
on a selection wedge that we now bound formally.

\paragraph{Selection into the administrative channel.}
Cases enter the administrative channel rather than the courts when the
patient (or a SUS social worker) requests the medication directly
through an SES/SP intake and the request is approved by a scientific
committee using cost-effectiveness criteria. The committee admits
clinically essential and sourceable cases; rejects expensive,
specialized, or uncertain ones. We do not treat the resulting selection
bias as ``ambiguous in sign''---the committee's first-stage probit on
pre-period item characteristics rejects a higher pre-period log
reference price ($-0.051$, $p<0.001$) and SUS-formulary status
($-0.265$, $p<0.001$): admin systematically over-represents items that
would have been priced lower under \emph{any} regime. The naive UTG
comparison therefore inflates the sanction effect.

\paragraph{Manski-Lee bounds.}
We bound the UTG coefficient under monotone selection by trimming
admin observations within item$\times$year$\times$PBU strata where the
admin sample exceeds the litigated count. The trimming proportion
$p_{\text{trim},s} = \max\{0,\,(n_{\text{adm},s}-n_{\text{lit},s})/n_{\text{adm},s}\}$
is the share of admin observations within stratum $s$ that, under the
Lee monotone-acceptance assumption, would not have appeared in the
litigated counterfactual. Trimming admin from the upper tail of the
within-stratum price distribution yields the lower bound on the
admin-minus-lit log-price gap; trimming from the lower tail yields the
upper bound. Average trimming intensity is \BPutgLeeTrimMean{} of admin
observations within strata where trimming applies; \BPutgLeeStrataExcess{}
of \BPutgLeeStrataTotal{} strata trigger trimming. Section~\ref{sec:utg_bounds}
reports the bounded UTG.

\paragraph{Heckman parametric correction (sensitivity).}
A parametric Heckman correction with the same first-stage probit and
item, year, and PBU FE in the second stage produces a second-stage
Admin coefficient of $\BPutgHeckmanCoef$ (SE $\BPutgHeckmanSE$, sample
size $\BPutgFirstStageN$). The inverse Mills ratio is collinear with
the second-stage fixed effects; the resulting point estimate is
uninformative. We treat the Heckman as a sensitivity diagnostic, not as
a competing point estimate, and rely on the Lee bound as the primary
selection-corrected object.

\paragraph{Inference.}
Standard errors are clustered at the PBU level. With approximately
\BPpbuCount{} clusters, asymptotic CR-cluster SEs may be liberal in the
tightest specifications. We re-do inference for the UTG coefficients
under a Rademacher wild cluster bootstrap ($B=\BPutgBoottestB$). The
bootstrap $p$-value on the preferred-FE Admin coefficient is
\BPutgBoottestPval; the corresponding 95\% CI in log-points is
$[\BPutgBoottestCIlow, \BPutgBoottestCIhigh]$. The naive-UTG result
survives wild-cluster inference; the within-item$\times$year-month
specification is borderline ($p=\BPutgBoottestPvalIYM$), as expected
given the thin within-cell sample documented below.

\paragraph{Within-cell representativeness.}
Item fixed effects absorb time-invariant item characteristics; item
$\times$ year-month FE absorb time-varying within-item composition.
The tightest specification therefore identifies off cells in which
both administrative and litigated purchases occur in the same
item$\times$year-month. Of \BPutgCellsItemYM{} such cells,
\BPbothCellN{} (about 11\%) carry within-cell variation; the remaining
\BPbothCellNSingleton{} are absorbed as singletons under the FE design.
Both-types cells are not a representative subsample (Online Appendix
Table~\ref{tab:both_types_cells}): they tilt toward higher-volume,
lower-priced, more-frequent SUS-formulary items than singleton cells.
The bounded UTG estimate therefore speaks to the more liquid items in
the urgent panel, not to the universe of court-mandated procurement.
We discuss this restriction explicitly in the welfare bound calculation
in Section~\ref{sec:results}.

\subsection{Specifications}

We estimate the urgent-vs-ordinary premium as
\begin{equation}
\label{eq:baseline}
y_{i,g,t} = \beta \cdot \text{Urgent}_{i,g,t} + \gamma_g + \lambda_t + \delta_b + \varepsilon_{i,g,t},
\end{equation}
where $y_{i,g,t}$ is the outcome (log negotiated price, log reference
price, log number of firms, success indicator) for purchase order $i$
of item $g$ at time $t$; $\gamma_g$, $\lambda_t$, $\delta_b$ are item,
time, and PBU fixed effects. We report four specifications with
progressively richer fixed effects.

We estimate the under-the-gun coefficient as
\begin{equation}
\label{eq:utg}
\ln(\text{Price}_{i,g,t}) = \beta \cdot \text{Admin}_{i,g,t} + \gamma_g + \lambda_t + \delta_b + \varepsilon_{i,g,t},
\end{equation}
on the urgent-only subsample (admin or litigated). We report (i) the
naive coefficient under the preferred FE; (ii) the Manski-Lee bounded
coefficient; (iii) the within firm-buyer-item triple coefficient that
isolates same-firm pricing from supplier composition.

We do not estimate a structural counterfactual price under hypothetical
sanction removal. That object is unidentified in our setting because
the administrative subsample is selected on item economics. What we
provide is a bounded descriptive decomposition of where the cost
margin between regimes comes from.
